The counterexample in Theorem~\ref{thm:s6-counterexample} changes the
six-element problem from a converse problem into a classification problem.
The accompanying laboratory therefore separates three levels:

\begin{enumerate}
\item paper-grade positive examples, such as the double-coset construction
      above;
\item paper-grade or nearly paper-grade obstruction families;
\item finite audit data whose current certificates are still too bulky for
      the main text.
\end{enumerate}

The current audit of the non-series-parallel six-element types satisfying
the necessary divisibility condition gives the following status.  Here
``unresolved'' means unresolved inside the current finite laboratory
pipeline; it does not mean that every negative row has a short
independently checkable paper certificate.
\[
\begin{array}{lr}
\text{non-series-parallel divisible candidates} & 57,\\
\text{non-series-parallel tilers} & 19,\\
\text{obstruction-side candidates classified negative by the audit} & 38,\\
\text{audit-unclassified survivors} & 0,\\
\text{negative audit rows still lacking compact independent certificates} & 5.
\end{array}
\]
This table is included as audit-level laboratory status data, not as a
theorem of this paper.  We do not promote the $19/38$ split to a theorem
because five obstruction-side rows do not yet have short independently
checkable certificates in the text or supplement.

\subsection*{Positive structure}

The positive rows in the current six-element audit include multiplier sets
with subgroup, coset, and double-coset structure.
Theorem~\ref{thm:s6-counterexample} is the cleanest instance:
the multiplier set is
\[
        A=HK\sqcup H\tau K,
\]
with $H\cong D_4$ and $K\cong S_3$.  In the verified witness, the
right-hand symmetry permutes a natural perfect matching in the cover graph
of the poset, while the left-hand symmetry organizes the target labels.
This suggests, but does not prove, that the positive side has a
group-theoretic organization beyond raw exact-cover search.

\subsection*{Obstruction grammar}

On the negative side, many obstruction-side candidates fall into reusable
families.  The current grammar includes:
\[
\begin{gathered}
\text{terminal-block divisibility},\qquad
\text{endpoint-pair obstructions},\\
\text{localized }(3,3,2)\text{ endpoint obstructions},\qquad
\text{positional gcd obstructions},\\
\text{single-position balance},\qquad
\text{ordered-pair balance},\\
\text{mod-}3\text{ peak/threshold obstructions in low partition modules}.
\end{gathered}
\]
These families are valuable because they explain failures by small
statistics of words, terminal blocks, or low-dimensional permutation
modules, rather than by a raw exact-cover search.
The language of low partition modules is the same representation-theoretic
environment behind partition-transitivity tests for sets of permutations
\cite{MartinSagan2006,FangXia2026}.

\subsection*{The five hard obstruction cases}

Five obstruction-side candidates remain the main compression problem:
\[
        001,\qquad 013,\qquad 017,\qquad 019,\qquad 020.
\]
These are the canonical IDs used in
\texttt{certificates/s6\_classification\_status.json}.  The canonical
ordering rule is part of the laboratory enumeration data rather than the
present paper text, so the IDs should be read as stable project-local
handles.

The current status table records them as defect-one residual packing
obstructions.  The numerical rows below are taken from
\texttt{s6\_classification\_status.json}; they are not derived from the
legacy Colab/RoundingSat timeout logs mentioned in the computational
appendix.  After forcing the identity translate, the audit computation
reports that the residual packing graph has maximum clique one below the
size required for a tiling.  In the table below, ``req.'' is the number of
tiles required for a full tiling,
``target'' is the residual clique size after the identity translate is
fixed, and ``conf.'' is the number of binary incompatibility constraints in
the residual graph:
\[
\begin{array}{c|rrrrrr}
\text{cand.} & |L(P)| & \text{req.} &
\text{target} & \text{max} & \text{verts.} & \text{conf.}\\
\hline
001 & 48 & 15 & 14 & 13 & 488 & 42850\\
013 & 24 & 30 & 29 & 28 & 313 & 6554\\
017 & 24 & 30 & 29 & 28 & 314 & 6498\\
019 & 24 & 30 & 29 & 28 & 314 & 6498\\
020 & 24 & 30 & 29 & 28 & 313 & 6554
\end{array}
\]
Equivalently, in the audit data, each maximum near-packing leaves a residual
set of exactly one tile-size, but the residual cannot be completed by another
translate.
The laboratory also sees these cases as integer holes in ordered-triple
marginals: the linear relaxation is feasible, but the integer decomposition
is infeasible.

The current paper does not promote these five rows to standalone theorems.
A satisfactory upgrade would be one of the following:

\begin{itemize}
\item a compact coloring or clique-cover certificate for the residual
      packing graph;
\item an orbit-quotient certificate under the stabilizer of the identity
      translate;
\item a residual-profile lemma showing that every optimal near-packing
      leaves a set of tile cardinality that is never a translate of $L(P)$;
\item a proof-logged pseudo-Boolean or SAT certificate with an independent
      verifier.
\end{itemize}

Until such certificates are extracted, the safe statement is that the
current finite audit records a $19/38$ laboratory status split among the
$57$ divisible non-series-parallel candidates, with five negative rows
still awaiting compact independent certificates.  The theorem-level
conclusion of this paper is the explicit failure of the series-parallel
converse.
